\documentclass[11pt,a4paper]{article}
\usepackage[margin=2.0cm]{geometry}
\usepackage[T1]{fontenc}
\usepackage{lmodern}
\usepackage{amsmath}
\usepackage{booktabs}
\usepackage{tabularx}
\usepackage{xcolor}
\usepackage{listings}
\usepackage[hidelinks]{hyperref}

\emergencystretch=2em
\definecolor{codebg}{gray}{0.95}
\lstset{language=C++, basicstyle=\ttfamily\footnotesize, keywordstyle=\color{blue!70!black},
  commentstyle=\color{green!50!black}\itshape, backgroundcolor=\color{codebg},
  frame=single, rulecolor=\color{black!20}, breaklines=true, showstringspaces=false,
  tabsize=2, columns=flexible}
\newcommand{\code}[1]{\texttt{#1}}

\title{ATmega32A Module \\ \Large Interrupts}
\author{Ali Sahafi}
\date{}

\begin{document}
\maketitle
\raggedright
\setlength{\parindent}{0pt}\setlength{\parskip}{4pt}

\section{Theory}

\paragraph{Polling or interrupt?} With \emph{polling} the program asks again
and again: ``Is the ADC finished? Has the timer overflowed? Has a byte
arrived?'' The CPU spends its time waiting. With an \emph{interrupt} the
peripheral tells the CPU itself when something happens. Until then the CPU is
free to do other work.

\paragraph{What happens at an interrupt.}
\begin{enumerate}
  \item An event happens (button edge, timer compare match, ADC finished, byte received\dots)
        and the peripheral sets its \emph{flag}.
  \item If this interrupt is enabled \emph{and} the global interrupt switch is on
        (\code{sei()}), the CPU finishes the current instruction.
  \item The CPU saves where it was and jumps to the \emph{Interrupt Service Routine} (ISR)
        of this event. While the ISR runs, other interrupts wait.
  \item At the end of the ISR the CPU goes back to the main program, exactly where it stopped.
\end{enumerate}

\paragraph{Vector table and priority.} The ATmega32A has 21 interrupt sources.
Each one has a fixed address in the \emph{interrupt vector table}. If two
interrupts come at the same time, the one with the lower address goes first
(\code{INT0} has the highest priority).

\paragraph{Rules for interrupt functions.}
\begin{itemize}
  \item Keep them \textbf{short}: no \code{\_delay\_ms()}, no long loops. Other interrupts wait while one runs.
  \item A variable that the interrupt function changes and \code{main()} reads must be
        \code{\textbf{volatile}}, for example \code{volatile uint8\_t count = 0;}
        Without \code{volatile} the compiler may never read the new value in \code{main()}.
\end{itemize}

\section{How to use the drivers}

In the drivers, your interrupt function is a \emph{normal function}. You give
it to the driver, and the driver writes the ISR for you, enables the interrupt
and calls \code{sei()}. All interrupt functions of the course drivers:

\begin{center}
\small
\begin{tabularx}{\textwidth}{@{}l l >{\raggedright\arraybackslash}X@{}}
\toprule
\textbf{Driver function} & \textbf{ISR (vector)} & \textbf{Your function is called \dots} \\
\midrule
\code{ExternalInterrupt.enable(INT0, when, f)} & \code{INT0\_vect} & on an edge of \code{PD2} (button S11). Also \code{INT1}, \code{INT2} \\
\code{Timer0.onOverflow(f)}                   & \code{TIMER0\_OVF\_vect}  & when Timer0 overflows \\
\code{Timer0.onCompareMatch(f)}               & \code{TIMER0\_COMP\_vect} & on a compare match (Timer1/Timer2: the same) \\
\code{ADC.onComplete(f)}                      & \code{ADC\_vect}          & when a conversion is finished: \code{f(value)}, value 0--1023.
                                                                           Start a conversion with \code{ADC.start(channel)} \\
\code{UART.onReceive(f)}                      & \code{USART\_RXC\_vect}   & when a byte has arrived: \code{f(byte)} \\
\midrule
\code{sei()} / \code{cli()} & & switch \emph{all} interrupts on / off (global switch) \\
\bottomrule
\end{tabularx}
\end{center}

Pass \code{nullptr} instead of a function to switch an interrupt off again
(\code{ExternalInterrupt.disable(INT0)} for the external interrupts).

\subsection*{External interrupts}

\begin{lstlisting}
ExternalInterrupt.enable(pin, when, function);              // no debounce
ExternalInterrupt.enable(pin, when, function, debounceMs);  // for a push button
ExternalInterrupt.disable(pin);
\end{lstlisting}

\begin{center}
\small
\begin{tabularx}{\textwidth}{@{}l l >{\raggedright\arraybackslash}X@{}}
\toprule
\textbf{Argument} & \textbf{Value} & \textbf{Meaning} \\
\midrule
\code{pin} & \code{INT0} & pin \code{PD2}: push button \textbf{S11} on the board \\
           & \code{INT1} & pin \code{PD3} (free pin on the header) \\
           & \code{INT2} & pin \code{PB2}: this is LED D2 on the board, so you can not use both \\
\midrule
\code{when} & \code{INT\_FALLING}   & HIGH $\rightarrow$ LOW: the button is \emph{pressed} \\
            & \code{INT\_RISING}    & LOW $\rightarrow$ HIGH: the button is \emph{released} \\
            & \code{INT\_ANY\_EDGE} & both edges (\code{INT0}, \code{INT1} only) \\
            & \code{INT\_LOW\_LEVEL} & again and again, as long as the pin is LOW (\code{INT0}, \code{INT1} only) \\
\midrule
\code{debounceMs} & \code{0} (default) & no debounce: every edge calls your function \\
                  & \code{1}--\code{255} & debounce time in ms, for example \code{20} for a push button \\
\bottomrule
\end{tabularx}
\end{center}

\code{enable()} makes the pin an input with pull-up. A wrong pin or edge is a
\textbf{compile error}, for example \code{ExternalInterrupt.enable(PD2, ...)}
(you must write \code{INT0}, not \code{PD2}).

\paragraph{Button bounce.} The push buttons on the board have no hardware
debouncer. When the metal contacts close, they bounce a few times, so one press
can make \emph{several} edges. We tested it on the board: 5 presses of S11 gave
the counter 7 without debounce. With \code{debounceMs = 20}, 10 presses gave
exactly 10. With a debounce time the driver waits after the first edge, forgets
the extra edges, and calls your function only if the pin is still at the new
level (for \code{INT\_FALLING}: still pressed). During this wait all other
interrupts must wait too, so use it only for buttons, not for fast signals.

\subsection*{Driver functions and registers}

\begin{center}
\small
{\footnotesize\begin{tabular}{@{}ll@{}}
\toprule
\textbf{Driver function} & \textbf{Registers} \\
\midrule
\code{ExternalInterrupt.enable(INT0, INT\_RISING, f)} & \code{DDRD \&= \textasciitilde(1<{}<PD2); PORTD |= (1<{}<PD2);} \\
            & \code{MCUCR |= (1<{}<ISC01) | (1<{}<ISC00);} \quad (rising edge) \\
            & \code{GIFR = (1<{}<INTF0); GICR |= (1<{}<INT0); sei();} \\
\code{ExternalInterrupt.disable(INT0)} & \code{GICR \&= \textasciitilde(1<{}<INT0);} \\
\code{Timer0.onCompareMatch(f)} & \code{TIMSK |= (1<{}<OCIE0); sei();} \\
\code{Timer0.onOverflow(f)}     & \code{TIMSK |= (1<{}<TOIE0); sei();} \\
\code{ADC.onComplete(f)}        & \code{ADCSRA |= (1<{}<ADIE); sei();} \\
\code{ADC.start(0)}             & \code{ADMUX = (ADMUX \& 0xE0) | 0; ADCSRA |= (1<{}<ADSC);} \\
\code{UART.onReceive(f)}        & \code{UCSRB |= (1<{}<RXCIE); sei();} \\
\bottomrule
\end{tabular}}
\end{center}

\newpage
\section{Example: up counter with button S11}

Every time button S11 is released, the interrupt adds 1 to the counter on the
LEDs. The main loop is empty. Try it and watch the bounce: sometimes the
counter jumps by 2 or 3. Then add a debounce time of 20\,ms as the fourth
argument of \code{enable()} and try again.

\paragraph{Build and upload.} \code{make interrupt}

\lstinputlisting{example.cpp}

\newpage
\section{Laboratory Tasks}\vspace{-4pt}
\setlength{\parskip}{2pt}\setlength{\itemsep}{0pt}

For each task, make a new \code{.cpp} file in the main project folder (for
example \code{Task1.cpp}). Start the file with \code{\#define F\_CPU 8000000UL}
and include the drivers you need. Upload it with \code{make SRC=Task1.cpp}.
In all tasks the main loop \code{while (true) \{ \}} stays \textbf{empty}: all
the work is done in interrupt functions.

\subsection*{Task 1: Real-Time Clock (Timer0 Interrupt)}
Count 0 to 9 on the seven-segment display, one step every \emph{exactly} 1
second (external 8\,MHz crystal). Use \code{Timer0} in \code{TIMER\_CTC} mode
with prescaler 8 and a compare match interrupt:
\begin{enumerate}
  \item With prescaler 8, one count takes $1\,\mu$s. Choose the compare value so that
        there is one compare match every $200\,\mu$s:
        \code{Timer0.setCompare(...)}.
  \item Give your tick function to \code{Timer0.onCompareMatch()}.
  \item In the tick function, count the ticks. How many ticks are 1 second? Then show
        the next digit with \code{SevenSeg(digit)}.
\end{enumerate}
Use \code{static} or \code{volatile} variables for the counters.

\subsection*{Task 2: ADC Interrupt}
Show the potentiometer on the LEDs, using the ADC interrupt instead of waiting:
\begin{enumerate}
  \item \code{ADC.enable()}, give your function to \code{ADC.onComplete()}, then start the
        first conversion with \code{ADC.start(0)}.
  \item Your function gets the value 0--1023. Show the upper 8 bits on the LEDs
        (\code{value >{}> 2}, active-low), then start the next conversion with
        \code{ADC.start(0)}.
\end{enumerate}

\subsection*{Task 3: UART Receive Interrupt}
Show every byte that the PC sends on the LEDs on \code{PORTB}. Give your
function to \code{UART.onReceive()}: it gets the received byte. Compare with
Lesson~6, Task~1.2: where does the program wait now?

\medskip\emph{Assignment 1 (the ultrasound sensor) is described on the lecture slides.}

\subsection*{Assignment 2: LED Brightness from the PC}
Receive bytes from the PC with \code{UART.onReceive()} and use each byte as the
brightness of LED D3: \code{Timer0} in \code{TIMER\_FAST\_PWM} mode with
\code{PWM\_INVERTING}, and \code{Timer0.setCompare(byte)} in your receive function.

\subsection*{Assignment 3: Measurements on Command}
The PC sends a command of \textbf{two bytes}: the first byte is the
\emph{number} of measurements, the second byte is the \emph{time} between them
in milliseconds. The board then measures the potentiometer this number of
times and sends every value back with \code{UART.println()}.
\begin{itemize}
  \item Receive the two bytes with \code{UART.onReceive()}. Store them in
        \code{volatile} variables and set a \code{volatile} flag when both bytes
        have arrived.
  \item In \code{main()}, when the flag is set, do the measurements. (Here the main loop
        is not empty: long work does not belong in an interrupt function.)
\end{itemize}

\end{document}
